\vkthmsync{8}
\paragraph{\trQuasiequivalence。}\label{sub:cunital-quasi-equivalences}
我们还需要一点更多的记号与术语：给定两个 \trCunital{} \trACategory{} $\scr{A}$
与 $\scr{B}$，记 $\mathit{fun}(\scr{A},\scr{B})$ 为
$\mathit{nu}\text{-}\mathit{fun}(\scr{A},\scr{B})$ 中由 \trCunital{} \trAFunctor{}组成的%
\trFull{}子\trCategory。一个 \trCunital{} \trAFunctor{} $\scr{F}\colon\scr{A}\to\scr{B}$
称为一个\textbf{\trQuasiequivalence}（quasi-equivalence），如果其底层的%
\trCohomology{}层面\trFunctor{}是一个\trEquivalence。

设 $\scr{A}$ 是一个 \trCunital{} \trACategory，而 $\tilde{\scr{A}}\subset\scr{A}$
是一个\trFull{} \trAinfty-子\trCategory{}，其包含是一个\trQuasiequivalence。

\begin{lemma}\label{thm:restriction-quasi-equivalence}
对任意 \trCunital{} \trACategory{} $\scr{B}$，把 \trCunital{} \trAFunctor{}及其预%
\trNaturalTransformation{}加以限制，便给出一个\trQuasiequivalence{}%
$\scr{Q}=\mathit{fun}(\scr{A},\scr{B})\to\tilde{\scr{Q}}
=\mathit{fun}(\tilde{\scr{A}},\scr{B})$。
\end{lemma}
\begin{proof}
任取两个 \trCunital{} \trAFunctor{} $\scr{G}_k$（$k=0,1$）以及它们在
$\tilde{\scr{A}}$ 上的限制 $\tilde{\scr{G}}_k$。预\trNaturalTransformation{}上的限制映射
\begin{equation}\label{eq:restriction-map-pre-natural}
\hom_{\scr{Q}}(\scr{G}_0,\scr{G}_1)
\longrightarrow\hom_{\tilde{\scr{Q}}}(\tilde{\scr{G}}_0,\tilde{\scr{G}}_1)
\end{equation}
显然与诸长度\trFiltration{}\trCompatible，因而诱导出相伴\trSpectralSequence{}之间的一个映射。在 $E_2$
层面上，它归结为相应\trCohomology{}\trFunctor{}的\trHochschildCohomology{}上的限制
映射，而我们刚刚证明了后者是一个\trIsomorphism。于是
\eqref{eq:restriction-map-pre-natural} 是一个\trQuasiisomorphism，也就是说，这个限制是%
\trCohomology\trFullyFaithful{}的。最后，给定任意一个 \trCunital{} \trAFunctor{}%
$\tilde{\scr{G}}\colon\tilde{\scr{A}}\to\scr{B}$，可以把它底层的\trCohomology{}\trFunctor{}%
唯一地扩张成一个 $G\colon H(\scr{A})\to H(\scr{B})$；于是\cref{thm:relative-functor-extension}
使我们得以构造一个 $\scr{G}\colon\scr{A}\to\scr{B}$，它扩张 $\tilde{\scr{G}}$ 且满足
$H(\scr{G})=G$，因而是\trCunital{}的。这证明了该限制在\trObject{}上是\trSurjective{}的，
从而是一个\trQuasiequivalence。
\end{proof}

特别地，这一扩张程序表明，存在一个 \trCunital{} \trAFunctor{}%
$\scr{P}\colon\scr{A}\to\tilde{\scr{A}}$ 使得
$\scr{P}|_{\tilde{\scr{A}}}=\mathrm{Id}_{\tilde{\scr{A}}}$。若以
$\scr{K}\colon\tilde{\scr{A}}\to\scr{A}$ 记该包含，则
$\scr{P}\circ\scr{K}=\mathrm{Id}_{\tilde{\scr{A}}}$。$\scr{P}$ 是一个%
\trQuasiequivalence，特别地是\trCohomology\trFullyFaithful{}的，因而由%
\cref{thm:left-composition-cff}，$\scr{L}_{\scr{P}}$ 亦然。由于
$\scr{L}_{\scr{P}}(\scr{K}\circ\scr{P})$ 等于 $\scr{L}_{\scr{P}}(\mathrm{Id}_{\scr{A}})$，
由此可知，作为 $H^0(\mathit{fun}(\scr{A},\scr{A}))$ 中的\trObject，$\scr{K}\circ\scr{P}$
与 $\mathrm{Id}_{\scr{A}}$\trIsomorphic。

\begin{theorem}\label{thm:quasi-equivalence-homotopy-inverse}
设 $\scr{F}\colon\scr{A}\to\scr{B}$ 是一个\trQuasiequivalence。那么存在一个%
\trQuasiequivalence{} $\scr{G}\colon\scr{B}\to\scr{A}$，使得在
$H^0(\mathit{fun}(\scr{A},\scr{A}))$ 中 $\scr{G}\circ\scr{F}\simeq\mathrm{Id}_{\scr{A}}$，
而在 $H^0(\mathit{fun}(\scr{B},\scr{B}))$ 中
$\scr{F}\circ\scr{G}\simeq\mathrm{Id}_{\scr{B}}$。
\end{theorem}
\begin{proof}
存在\trFull{} \trAinfty-子\trCategory{}%
$\tilde{\scr{A}}\subset\scr{A}$、$\tilde{\scr{B}}\subset\scr{B}$，使得诸包含是%
\trQuasiequivalence，并且 $\scr{F}$ 限制为一个从 $\tilde{\scr{A}}$ 到 $\tilde{\scr{B}}$
的\trQuasiisomorphism{} $\tilde{\scr{F}}$。正如我们已然确定的，这两个包含都可以在各自的%
\trFunctor{}\trCategory{}中于\trIsomorphism{}的意义下求逆。$\tilde{\scr{F}}$ 本身由%
\cref{thm:quasi-iso-inverse-homotopy}可在\trHomotopy{}意义下求逆，而由\cref{thm:homotopic-functors-isomorphic}，
在\trFunctor{}\trCategory{}中\trHomotopy{}蕴涵\trIsomorphism。
\end{proof}

\begin{lemma}\label{thm:right-composition-quasi-equivalence}
如果 $\scr{G}\colon\scr{A}\to\scr{B}$ 是一个\trQuasiequivalence，那么对任意 \trCunital{}
\trACategory{} $\scr{C}$，$\scr{R}_{\scr{G}}\colon\mathit{fun}(\scr{B},\scr{C})
\to\mathit{fun}(\scr{A},\scr{C})$ 亦然（关于 $\scr{L}_{\scr{G}}$ 的平行命题同样成立；
事实上，此时对\trNonunital{}的 $\scr{C}$ 以及诸\trCategory{}%
$\mathit{nu}\text{-}\mathit{fun}$ 也存在相应的版本）。
\end{lemma}
\begin{proof}
我们首先需要知道 $\scr{R}_{\scr{G}}$ 是\trCohomology\trFullyFaithful{}的。若 $\scr{G}$
是一个\trQuasiisomorphism，这是\cref{thm:left-composition-cff}的一部分内容；而若 $\scr{G}$
是一个\trFull{}子\trCategory{}的包含，则是\cref{thm:restriction-quasi-equivalence}的一部分
内容。与上文的讨论相平行，这两个特殊情形蕴涵一般结果。接下来，%
\cref{thm:quasi-equivalence-homotopy-inverse}给出一个 $\scr{F}\colon\scr{B}\to\scr{A}$
使得 $\scr{F}\circ\scr{G}\simeq\mathrm{Id}_{\scr{B}}$。由\cref{thm:left-composition-cunital}，
对任意 $\scr{H}\in\Ob\mathit{fun}(\scr{A},\scr{C})$ 可得
$\scr{R}_{\scr{G}}(\scr{H}\circ\scr{F})
=\scr{L}_{\scr{H}}(\scr{F}\circ\scr{G})
\simeq\scr{L}_{\scr{H}}(\mathrm{Id}_{\scr{B}})=\scr{H}$，这蕴涵 $\scr{R}_{\scr{G}}$
在\trObject{}的\trIsomorphism{}类上是一个\trSurjective。关于 $\scr{L}_{\scr{G}}$
的论证类似，只是稍微简单一点。
\end{proof}

\begin{remark}\label{thm:strictly-unital-version-history}
\cref{thm:quasi-equivalence-homotopy-inverse}的\trStrictlyUnital{}版本看来最早是由%
\trKontsevich{}陈述的（未发表）。它出现于 [55, Theorem 8.6] 与 [95, Theorem 9.2.04]
中。此处讨论的\trCunital{}版本也可以在 [97, Theorem 8.8] 中找到。
\end{remark}
